\section{The three finite-coefficient local conditions at a large prime}

% BEGIN T80B1
\paragraph{T80B1. The gap between the letter block and the other jumps.}
Use the native numerical datum at an unramified field $K/\mathbf Q_\ell$.
At embedding $i\in\mathbf Z/n$ the distinguished jumps form the block
$[w_i,w_i+t-1]$. The other jumps, with multiplicities, form $U_i$.
Write $h_i$ for the height, $w'_i=h_i-w_i$, and put
\[
 s_i=\left\lfloor\sum_{j\ge1}\ell^{-j}h_{i+j}\right\rfloor,
 \qquad d_i=\#\{u\in U_i:u<w_i\},\qquad d=\sum_i d_i .
\]
The series converges because the heights are periodic. Multiplying the
preceding series by $\ell$ and taking integer parts gives
$h_i+s_i+1\le\ell(s_{i-1}+1)$. Put $\eta_i$ for the width of the
short end of the height interval. The three numerical hypotheses impose
$\eta_i+4s_i+6<\ell$, distance $>s_i+2$ between every low jump and
every block entry, and exclusion of cloud residues throughout
$[-2s_i-2,\eta_i+2s_i+3]$, measured from $0$ at a good letter and
from $h_i$ at a boundary letter. Low jumps and the block lie in
$[1,\eta_i]$ at a good letter, in $[h_i-\eta_i,h_i-1]$ at a boundary
letter. Consequently no $u\in U_i$ is congruent modulo $\ell$ to any
integer in
\[
 [w_i-s_i-2,w_i+t+1+s_i].                                      \tag{T80.1}
\]
Indeed for cloud entries this interval lies in the indicated exclusion
interval in the appropriate coordinate. For low entries a congruence
in this interval must be an equality, since the integer difference has
absolute value $<\ell$. Projection onto the nearest block entry then
violates the distance hypothesis. The same argument at the short end
excludes residues of $U_i$ in $[-s_i,0]$ at good letters and in
$[h_i,h_i+s_i+1]$ at boundary letters. The cloud misses
$[0,\eta_i]$ or $[h_i-\eta_i,h_i]$, respectively.
% END T80B1

% BEGIN T80B2
\paragraph{T80B2. Which residual characters can occur.}
Assume the native residual realization $S.\mathrm{Real}$, and fix its
model over $k=O/\lambda$. Thus the twisted lattice
$T=T_0\otimes\chi_1^{-1}$ reduces to its specified object $A$, the
model functor is exact on the indicated subobjects and quotients, and it
supplies base models with their prescribed Sen multiplicities, character
models and polynomial identities. Write $v$ for the reduced Wach
parameter and $\nabla=(\gamma-1)/v$.
For a model lattice $N_i$, let $m_N(i,y)$ be the multiplicity of the
generalized $y$-eigenspace of $\nabla$ on $N_i/vN_i$. In the source
normalization multiplication by $v^a$ changes the eigenvalue $y$ to
$y-a$. This follows from $\gamma(v)/v=1-v+O(v^2)$ and semilinearity.
Generalized eigenspace dimensions add along finite invariant filtrations.
In particular on $v^aN_i/v^{a+b}N_i$ the zero multiplicity equals
$\sum_{j=0}^{b-1}m_N(i,a+j)$.

Here are the lattice comparisons needed below. The two height inclusions
are $v^{w'_i}N_i\subset P_i\subset v^{-w_i}N_i$, with
$P_i=k[[v]]\varphi_i N_{i-1}$. For any two such families the least
nonnegative exponents $a_i$ satisfying $v^{a_i}N_i\subset N'_i$
obey $\ell a_{i-1}\le a_i+h_i$. Indeed $v^{w_i}\varphi_i x\in N_i$
for $x\in N_{i-1}$, so $v^{a_i+h_i}\varphi_i x\in P'_i$. Express
$x$ in a lattice basis of $N'_{i-1}$. The Frobenius images of this
basis are a basis of $P'_i$, and Frobenius multiplies the order of each
coefficient by $\ell$. Iterating the inequality and using boundedness
on the finite cycle gives $a_i\le s_i$. Interchanging the two families
sandwiches each between the shifts $v^{s_i}N'_i$ and $v^{-s_i}N'_i$.
Its Sen weights therefore differ from some weight of $N'_i$ by an
integer of absolute value at most $s_i$. To see this, embed the
subquotient $N_i/vN_i$ as a subquotient of
$v^{-s_i}N'_i/v^{s_i+1}N'_i$ and filter by powers of $v$.

Take a subquotient of $A$ with any model. The model induced from the
standard lattice by intersection with the subspace and image in the
quotient has weights among $U_i-w_i$: both exact sequences of lattices
have free quotients, so remain exact modulo $v$. An arbitrary model
has weights among $(U_i\cup[w_i,w_i+t-1])-w_i$ by its realization
as a subquotient of a base model, using the same argument. A weight
from the block would have to be within $s_i$ of a residue from
$U_i-w_i$ by the sandwich just proved. This contradicts (T80.1).
In particular $A$ has no subquotient $k(\bar\epsilon^{-e})$ for
$-1\le e<t$. Its prescribed line lattice has Sen weight $e$. It has the two
height inclusions since $-w_i\le e\le w'_i$ and its Frobenius pull
is $v^eN_{e,i}$. Thus it may be compared even for $e=-1$, where
no model of that line is assumed. The induced subquotient lattice would
produce a residue of $U_i-w_i$ within $s_i$ of $e$, forbidden by
(T80.1).

Every model of $k(\bar\epsilon^{-e})$ for $0\le e<t$ equals the
prescribed line lattice. Indeed any such lattice is $v^{a_i}N_{e,i}$
at $i$, has weight $e-a_i$, and $|a_i|\le s_i$. Frobenius on the
standard line has lattice image $v^eN_{e,i}$. Its two height inclusions give
\[
 -w_i+a_i\le\ell a_{i-1}+e\le w'_i+a_i .
\]
At a good letter $w_i+e\le\eta_i$ forces $a_{i-1}\ge0$, and at a
boundary letter $w'_i-e\le\eta_i$ forces $a_{i-1}\le0$.
Take a composition series of the base object passing through the given
line subquotient, inducing models on its factors. Such a series is
obtained by refining kernel and image chains. Residual semisimplification
gives precisely the factors of copies of $A$ and equally many copies of
each indicated line. The two lists are disjoint by the preceding exclusion,
and the powers in the line list are distinct by the cyclotomic order in
the realization. On the integer interval $[-s_i,t-1+s_i]$, the factors
of $A$ contribute no weights, by (T80.1). The remaining weights, counted
with multiplicities, must therefore be the original list of $e$'s by
the base-model Sen formula. The interval has length $<\ell$, so taking
the ordinary integer sum gives $\sum a_i=0$ over the line factors.
They all have the same weak sign fixed by the letter at $i+1$. Each shift
is zero, in particular the shift of the given line.
% END T80B2

% BEGIN T80B3
\paragraph{T80B3. From a model extension to polar parts.}
At $X=k$, admission for every sufficiently large exponent supplies an
admitted extension at exponent one: map $(x,y\bmod\lambda^a)$ onto
$(x,y\bmod\lambda)$. This is equivariant because the upper entry is
killed by $\lambda$, and a quotient of a lattice quotient is again a
lattice quotient, by taking the preimage of the new kernel.
After twisting this admitted residue extension back by $\chi_1^{-1}$,
the realization gives a model of $0\to A\to E\to k\to0$.
Its quotient lattice is standard by T80B2. Lift the unit section at
each letter to $u_i$ in the middle lattice. In the kernel put
$b_i=\varphi_i u_{i-1}-u_i$ and $g_i=(\gamma_i-1)u_i$.
Changing the lifts adds the pair
$(\varphi_i y_{i-1}-y_i,(\gamma_i-1)y_i)$. The relation is
$(\gamma_i-1)b_i=\varphi_i g_{i-1}-g_i$. A zero class of pairs gives
a section fixed by both operators. The fixed-vector comparison
in the model functor gives a Galois map from the trivial line taking
the model unit to this corrected lift. Its composite with the
projection is scalar on the trivial line and fixes the nonzero model
unit, so that scalar is one. This is a Galois splitting.
Addition and scalar multiplication of extensions respect these pairs:
use the diagonal on the quotient and addition or scalar multiplication on
the kernel. Thus a linear relation modulo pair coboundaries gives a
relation modulo Galois coboundaries.

In the induced kernel lattice $N$, $g_i\in vN_i$. At good letters
$b_i\in v^{-w_i}N_i$. At boundary letters $b_i\in v^{-w'_i}P_i$
and there is an $x_i\in N_i$ with $g_i-(\gamma_i-1)x_i\in
v^{w'_i+2}N_i$. For the first boundary assertion, write
$v^{w'_i}u_i$ in the Frobenius pull of the middle lattice. That pull
is the sum of the pull of the kernel and the span of
$\varphi_i u_{i-1}$, since each preceding lattice vector splits as a
kernel vector plus an integral multiple of the preceding lift.
The coefficient of $\varphi_i u_{i-1}$ is $v^{w'_i}$ by projection
to the unit. For the second assertion the base polynomial condition
gives $f(0)\ne0$ and
$\nabla f(\nabla)N_{E,i}\subset v^{w'_i+1}N_{E,i}$.
The condition passes to intersections and images, since maps commute
with $\nabla$. Use $x_i=u_i-f(0)^{-1}f(\nabla)u_i$, which projects
to zero and has the required difference.

All kernel lattices can now be enlarged to a single greatest model of
$A$. Models are closed under sums: take the direct sum of two base
presentations, pull its lattice back to the direct sum of the two
subobjects, and push it forward by addition into $A$. The two
coordinate projections identify the resulting image with the sum. T80B2 bounds their sum inside $v^{-s_i}N_{0,i}$
at each letter, where $N_0$ is the standard model. Noetherianity and
finiteness of the cycle realize the sum by finitely many models. Write
$N$ for this greatest model. The conditions on $b_i,g_i,x_i$ persist
under enlargement.

For this $N$, write $m_i(y)=m_N(i,y)$ and $m^0_i$ for the standard
multiplicity. Let $q_i(y)$ be the generalized $y$-multiplicity on
$N_i/N_{0,i}$. Filtering $N_i/vN_{0,i}$ first by $N_{0,i}$, then
by $vN_i$, and using the shift by $v$, gives
\[
 m_i(y)+q_i(y+1)=m^0_i(y)+q_i(y).
\]
Since $N_i\subset v^{-s_i}N_{0,i}$, a nonzero $q_i(y)$ requires
$y-j=u-w_i\pmod\ell$ for some $1\le j\le s_i$ and $u\in U_i$.
In particular $q_i(0)=0$, and $m_i(a)=0$ for $0\le a\le s_i+1$
by T80B2 and (T80.1). At a good letter $q_i(-w_i)=0$ by the
short-end exclusion. Telescoping over $[-w_i,-1]$ gives
$\sum_{y=-w_i}^{-1}m_i(y)=d_i$: an actual low entry in
$[0,w_i-1]$ is counted once, and the cloud misses the interval.
At a boundary letter telescope instead from $0$ through $w'_i+s_i+1$.
The terminal $q_i(w'_i+s_i+2)$ and the terms $m_i(y)$ for
$w'_i<y\le w'_i+s_i+1$ vanish by the upper-end exclusion of T80B1.
On the standard lattice the count is the number of entries $u\ge w_i$,
since congruence for low entries in the short interval is equality,
and the cloud misses the short interval and its upper extension. Thus
$\sum_{y=0}^{w'_i}m_i(y)=r-d_i$.
% END T80B3

% BEGIN T80B4
\paragraph{T80B4. Division by $\gamma-1$ on the kernel of $\psi$.}
Use the cyclotomic embedding supplied by the realization. In the full
parameter $u$ it identifies $k((v))$ with the invariants of the
Teichm\"uller group $\Delta$ in $k((u))$ and multiplies orders by
$\ell-1$. Write $q=1+u$, $\Phi(u)=u^\ell$, and
$\gamma(q)=q^{1+\ell}$. There is a unique expansion
$z=\sum_{j=0}^{\ell-1}q^j\Phi(z_j)$: collect powers of $u$ modulo
$\ell$ and use the triangular binomial basis $1,q,\ldots,q^{\ell-1}$
of the polynomials of degree $<\ell$. Here $\Phi$ fixes $k$. The
coefficient projection $\psi(z)=z_0$ has
$\psi(\Phi(a)z)=a\psi(z)$, commutes with $\gamma$, and commutes with
$\Delta$. Indeed substitution $q\mapsto q^d$ for a unit exponent
permutes the nonzero residue classes $dj\bmod\ell$, multiplying each
summand by a Frobenius factor. Thus it restricts to the invariant ring.

For $2-\ell\le a\le1$, $\psi(v^ak[[v]])=k[[v]]$. In full
coordinates its coefficient of degree $m$ is
$\sum_{j=0}^{\ell-1}(-1)^j[z]_{\ell m+j}$, by the inverse binomial
change of basis. For an input of order $\ge(\ell-1)a$ this gives an
output of order $\ge 2-\ell$ and thus nonnegative order in the
invariant coordinate. Moreover $\psi(v)$ has nonzero constant term:
only the coefficient of degree $\ell-1$ can contribute, and $v$ has
exactly that order. Multiples of $v$ of the form
$v\Phi(b/\psi(v))$ therefore realize every integral output $b$.

For a Frobenius pull $P=k[[v]]\varphi N$ define $\Psi$ by
$\Psi(c\varphi x)=\psi(c)x$. This exists by writing vectors in a
Frobenius image basis, and the displayed law makes it basis-independent.
It commutes with $\gamma$ between consecutive letters. For
$2-\ell\le a\le0$ the map
\[
 \gamma-1:\ker\Psi\cap v^aP\longrightarrow
                  \ker\Psi\cap v^{a+1}P
\]
is bijective. To check the valuation exactly, a scalar $z\in\ker\psi$
lies in $v^a k[[v]]$ iff its full-coordinate image lies in
$u^{\ell a}k[[u]]$ for $2-\ell\le a\le1$. For $a\le0$ this follows
from $(\ell-1)(a-1)<\ell a\le(\ell-1)a$. For $a=1$ the
possible term of degree $\ell-1$ must vanish by the constant-coefficient
formula for $\psi$. In full coordinates, having order $\ge\ell a$ is
equivalent to every coefficient $z_j$ having order $\ge a$, by the
binomial basis change. In a basis of $N$, let $H$ be the matrix of its
semilinear $\gamma$ after embedding into $k((u))$. Then $H-I=O(u^2)$
because $(\gamma-1)N\subset vN$ and $\ell\ge5$. For the coefficient
$j\ne0$ of an element of the pull, $\gamma-1$ acts as
$y\mapsto q^j H\gamma(y)-y$; this follows from
$\gamma(q^j)=q^j\Phi(q^j)$. On every graded piece it raises order by
one with leading coefficient $j$, since $\gamma(u)/u=1+O(u^2)$
and $H-I=O(u^2)$. Successive coefficient solving therefore gives a
unique preimage of order $\ge a$ for each input of order $\ge a+1$.
The $j=0$ term vanishes in the kernel. Uniqueness makes the regrouped
solution $\Delta$-invariant whenever its target and matrix entries
descend. This proves the assertion on the original pull.
% END T80B4

% BEGIN T80B5
\paragraph{T80B5. The residue dimension of the rank-$r$ condition.}
The zero multiplicities from T80B3 make $\nabla$ bijective on
$N_i/v^{s_i+2}N_i$, by filtering it by the individual powers. For
each pair choose $y_i\in N_i$ solving
$\nabla y_i\equiv g_i/v\pmod{v^{s_i+2}N_i}$ and subtract its
coboundary. Call the resulting pair $(b'_i,g'_i)$. Then
$g'_i\in v^{s_i+3}N_i$. At good letters put $c_i=b'_i$.
At boundary letters put $c_i=b_i-\varphi_i y_{i-1}+x_i$ with $x_i$
as in T80B3. In either case $c_i-b'_i\in N_i$. The cocycle identity,
the height inclusions, and $\ell(s_{i-1}+3)-w_i\ge w'_i+2$ give
\[
 \begin{array}{ll}
 c_i\in v^{-w_i}N_i,\quad \nabla c_i\in N_i,
                                      &i\text{ good},\\
 c_i\in v^{-w'_i}P_i,\quad \nabla c_i\in v^{w'_i+1}N_i,
                                      &i\text{ boundary}.
 \end{array}                                                    \tag{T80.2}
\]
For example $(\gamma_i-1)c_i=\varphi_i g'_{i-1}-
(g_i-(\gamma_i-1)x_i)$ at a boundary letter.

On any independent family, recording the $c_i\bmod N_i$ for each
member and extending linearly is injective. Indeed if all polar parts
of a combination vanish, its $b'_i$ are integral. The cocycle identity
makes $\nabla b'_i\in v^{s_i+2}N_i$, so $b'_i\in v^{s_i+2}N_i$
by the same finite-window bijectivity. On the product of these last
lattices Frobenius raises the filtration by at least one because
$\ell(s_{i-1}+2)-w_i\ge s_i+3$. The convergent geometric series
solves $(\varphi-1)y=b'$ uniquely. Subtraction leaves a Frobenius-fixed
$g'-(\gamma-1)y$ in this contracting product, hence zero.

At a good letter the polar span in (T80.2) embeds in the kernel of
$\nabla$ on $v^{-w_i}N_i/N_i$. Its dimension is at most the
generalized zero multiplicity $\sum_{y=-w_i}^{-1}m_i(y)=d_i$.
At a boundary letter abbreviate $a=w'_i$, $N=N_i$, $P=P_i$,
$U=v^{-a}P$, and $L=v^{a+1}N$. Then $L\subset N\subset U$.
In the ambient space modulo $L$ let
$B=\{x\in U:\nabla x\in L\}/L$. We show $\dim B\le r$.
Choose a lattice basis $e_j$ of $N_{i-1}$ and let $J$ be the
$k$-span of $v^{-1}\varphi_i e_j$. We have
$U\subset J+\nabla U+L$. For $x\in U$ set
$y=\Psi(vx)\in N_{i-1}$ by T80B4 with exponent $1-a$.
There exists $x'\in U$ with
$(\gamma_i-1)x'=vx-\varphi_i y$, since the right side lies
in $\ker\Psi\cap v^{1-a}P$. Here $a\ge1$ and $a+7\le\ell$ by T80B1.
Write $y=y_1+vy_3$ with $y_1$ in the span of the $e_j$. Choose
$z\in N_{i-1}$ with $y_3-\nabla z\in v^{s_{i-1}+2}N_{i-1}$.
Commuting Frobenius with $\gamma$ gives
\[
 x=v^{-1}\varphi_i y_1+\nabla(x'+\varphi_i z)
    +v^{-1}\varphi_i\big(v(y_3-\nabla z)\big).
\]
The final term belongs to $L$ by $\ell(s_{i-1}+3)-w_i\ge a+2$.
Now $U/L$ has finite dimension. The map it induces by $\nabla$ into
the ambient quotient has kernel $B$; it need not preserve $U$.
The containment just proved and rank-nullity give
$\dim B\le\dim J\le r$.

On $N/L$ the base polynomial of T80B3 annihilates
$\nabla f(\nabla)$, with $f(0)\ne0$. Its generalized zero eigenspace
is therefore its ordinary kernel, of dimension
$\sum_{y=0}^a m_i(y)=r-d_i$. This subspace of $B$ maps to zero modulo
$N$, whereas all polar parts in (T80.2) lie in the image of $B$ modulo
$N$. Their dimension is consequently at most $d_i$. Adding the letter
bounds and applying injectivity proves that at most $d$ independent
admissible residue classes occur. In using the model, represent a class
by its cocycle and its triangular residue extension. Reduction identifies
$T\otimes k$ with $A$ and the extension with $A\oplus k$ with upper
entry that cocycle. An independent family in Galois cohomology remains
independent in the model-pair quotient by T80B3.
% END T80B5

% BEGIN T80B6
\paragraph{T80B6. Zero and one for the character lines.}
For the $c$th line, twisting by $\chi_1^{-1}$ gives the residual
character $\bar\epsilon^{-e}$ with $e=c-1$. Its lattice must be
$N_e$ by T80B2. First suppose $1\le e<t$. Put $J=w_i$ at a good
letter and $J=w'_i-e$ at a boundary letter. In both cases
$0\le J$, $e+J<\ell$, $b_i\in v^{-J}N_{e,i}$, and
$g_i\in vN_{e,i}$. Frobenius takes $N_{e,i-1}$ to a lattice
whose span is $v^eN_{e,i}$, so the identity gives $\nabla b_i\in N_{e,i}$.
The eigenvalues on $v^{-J}N_{e,i}/N_{e,i}$ are
$e+1,\ldots,e+J$, all nonzero. Thus $b_i$ is integral. On the
integral product Frobenius is contracting since $e\ge1$. Solve
$(\varphi-1)y=b$, then the remaining $g-(\gamma-1)y$ is a fixed
point in the contracting product. It is zero. The original extension
splits and its residue class vanishes.

For $e=0$ use $J=w_i$ or $J=w'_i$. The same polar argument works,
since $J\le\eta_i<\ell$. In the standard unit basis all $b_i$ are
integral and $g_i$ have zero constant coefficient. Subtract the constant
coefficients $c_i$ of $b_i$. Frobenius contracts the positive-order
product, so the preceding argument makes the pair equivalent to
$(c_i,0)$. A tuple of constants with $\sum_i c_i=0$ has the form
$c_i=y_{i-1}-y_i$ by cyclic summation, so it too represents zero.
Every two admitted classes have a linear combination with this sum
zero, giving residue dimension at most one. This works also when
$\ell$ divides the number of letters.
% END T80B6

% BEGIN T80B7
\paragraph{T80B7. Cohomology of the twisted lattice and its finite summand.}
Let $H=H^1(K,T)$ and $V=T[1/\lambda]$. The case $e=0$ of
T80B2 excludes a trivial subrepresentation of $\bar T$; the case
$e=-1$ excludes a cyclotomic quotient. Thus $H^0(K,\bar T)=0$.
The local Tate pairing identifies $H^2(K,\bar T)$ with the dual of
the invariants of the finite Tate dual, which likewise vanish.
Indeed a nonzero equivariant additive map to $\mu_\ell$ gives a
nonzero $k$-linear map to $k(\bar\epsilon)$ by the nondegenerate
trace pairing on $k/\mathbf F_\ell$: recover its value on $x$ by
pairing with every scalar multiple of $x$. Euler characteristic now gives
$\#H^1(K,\bar T)=\#\bar T^{[K:\mathbf Q_\ell]}$, so its
$k$-dimension is $nr$.

Write $T_a=T/\lambda^aT$. Induction on $a$ using the long exact
sequences gives $H^0(K,T_a)=H^2(K,T_a)=0$ and surjective maps
$H^1(K,T_{a+1})\to H^1(K,T_a)$ with kernel $H^1(K,\bar T)$.
In particular $\#H^1(K,T_a)=|k|^{anr}$. In degree one continuous
lattice cohomology is the inverse limit: compatible classes admit
compatible cocycles by lifting the correcting zero-cochains; a compatible
system of coboundaries has compatible primitives because $H^0(K,T_a)=0$.
Projection of the limit to $H^1(K,T_a)$ induces reduction modulo
$\lambda^a$: in the exact sequence for $\lambda^a:T\to T$, the
kernel is $\lambda^a H$, and the just-proved inverse-limit description
gives surjectivity. Multiplication by $\lambda$ on $H$ is injective
since $H^0(K,\bar T)=0$. Lift a basis in degree one at $a=1$.
Successive reduction and completeness span $H$ by these $nr$ vectors,
and a primitive relation would contradict their independence modulo
$\lambda$. Thus $H$ is free of rank $nr$.
Decomposing a finite $O$-module $X$ into cyclic modules now proves,
for the natural map,
\[
 H\otimes_O X\ \simeq\ H^1(K,T\otimes_O X).
\]

The Bloch--Kato local dimension formula for a crystalline representation
is $\dim H^1_f(K,V)=\dim H^0(K,V)+\sum_{i,a<0}m_V(i,a)$ in these
jump conventions. Rational invariants vanish, since an invariant vector
can be scaled into $T\setminus\lambda T$. Twisting subtracts $w_i$
from every entry of $U_i$. Hence the dimension is $d$.
The inverse image $H_f$ of $H^1_f(K,V)$ in $H$ is saturated, hence a
direct summand of rank $d$. For the rank assertion, a rational cocycle
is bounded on the compact Galois group and hence becomes integral after
one scaling; the same is true for a rational primitive. Thus
$H[1/\lambda]=H^1(K,V)$ and $H_f$ spans its designated rational
subspace.
% END T80B7

% BEGIN T80B8
\paragraph{T80B8. Admissibility over finite coefficient modules.}
For a lattice $W$ equal to $T_0$ or one of the indicated character
lattices, let $A_W(X)$ be the set of admitted cohomology classes for
the twist of $W$ by $\chi_1^{-1}$. At an exponent $a\ge1$ killing
$X$ by $\lambda^a$, admission concerns the triangular extension with
quotient $O/\lambda^a$ and with the original $\chi_1$ twist restored.
The underlying class of admissible representations is closed under sums
and subquotients: direct sums of exhibiting crystalline lattices add the
multiplicities and residual factors; a quotient composes the quotient
map and for a subobject one takes its preimage. The untwisted coefficient
objects $W\otimes X$ and scalar quotients $O/\lambda^a(\chi_1)$
are admitted: take the standard direct sum with all its prescribed
summands and pass to subquotients of copies of its lattice.

Consequently $A_W(X)$ is a submodule. For addition, pull two extensions
back on the diagonal quotient line and push out by addition of their
kernels. Pushing out handles arbitrary coefficient maps. For an injection
$X\hookrightarrow X'$ the smaller extension embeds in the larger one
at a common exponent, since $W$ is free. Thus membership also pulls back
along injections. A change of cocycle by a coboundary changes the
splitting and preserves admission. Enlarging the exponent is pullback
along $O/\lambda^b\to O/\lambda^a$ and passage to a subobject of
the old extension times the new quotient line; decreasing it is passage
to the quotient. This proves that admission at one killing exponent is
equivalent to admission at every positive killing exponent and to the
eventual condition in the native definition.

For $T_0$, $H_f\otimes X\subset A_{T_0}(X)$ under T80B7.
Indeed an integral cocycle with finite rational class gives a crystalline
extension of the trivial character by $V$: its class becomes a
coboundary over $B_{\rm cris}$, so the middle term splits over that
ring and has full crystalline period dimension. Twist the extension by
$\chi_1$. Strictness of the de Rham filtration adds precisely the jump
$w_i$ to $U_i$ at each embedding. The lattice defined by the triangular
matrix reduces to an extension with the same residual factors, and
adjoining the remaining character summands gives exactly the prescribed
jump multiset and semisimplification. At any killing exponent reduce
this triangular lattice and project away those additional summands to
get the needed admitted extension for a cyclic coefficient. Maps
$O/\lambda^a\to X$ and addition give the asserted inclusion.

Choose a complement $C$ to $H_f$ in $H$. The admitted vectors lying in
$C\otimes k$ vanish: otherwise such a vector together with a basis of
$H_f\otimes k$ contradicts the bound $d$ of T80B5. For a nonzero
finite $X$ choose a simple submodule $k\hookrightarrow X$. By induction
on length the image in $C\otimes(X/k)$ of an admitted vector in
$C\otimes X$ vanishes. Lift it to $C\otimes k$ by exactness, and use
pullback of membership to see this lift vanishes too. Subtracting its
$H_f$ component from any admitted vector therefore proves
$A_{T_0}(X)=H_f\otimes X$.

For $2\le c\le t$ the zero bound of T80B6 similarly implies that the
only admitted class for $O(\chi_c)$ is zero over every $X$. In this
argument use exactness of $H^1$ at the middle of the cohomology sequence
for $k\hookrightarrow X\twoheadrightarrow X/k$ with the indicated
line twist: if the image is zero, subtract a lifted coboundary to obtain
a cocycle in the smaller module, where injective pullback gives admission.
This proof needs no freeness statement for the cohomology of these lines.
% END T80B8

% BEGIN T80B9
\paragraph{T80B9. The unramified line and the exact three conclusions.}
On the first character line the twisted action is trivial. Write
$\nu:G_K\to\mathbf Z_\ell$ for the unramified degree with
$\nu(F)=1$ for a Frobenius lift $F$. An unramified homomorphism
to finite $X$ has form $z(g)=\nu(g)\xi$. We verify its admission,
including the prescribed residual polynomial. Choose a positive
killing exponent $b=\lambda^a$, and a power $c=\ell^m$ divisible
by $b$, $m\ge1$. Set $\mu(g)=(1+c)^{\nu(g)}$ and
$w(g)=(\mu(g)-1)/c$. Then $w(g)\xi=\nu(g)\xi$ for odd $\ell$:
for integral nonnegative degrees expand the binomial modulo $c^2$
and extend by continuity. Both $\chi_1$ and $\chi_1\mu$ are
crystalline of the same jumps and agree modulo $\lambda$. Replacing
the first line by either one in the prescribed direct sum therefore
supplies admitted reductions modulo $c^2$ of both lines.

The triangular extension of $O/b$ by $X$ for $z$ is an explicit
subquotient. Inside $X\oplus(O/c^2)^2$, with actions
$\chi_1,\chi_1,\chi_1\mu$, take the image of
$(x,s,y)\mapsto(x,cs-y,y)$ for $s,y\in O$.
Map it onto $X\oplus O/b$ by $(x,s,y)\mapsto(x+s\xi,y\bmod b)$
on representatives. The first kernel is contained in the second:
it implies $x=0$, $c^2\mid y$, and $c\mid s$. The action on
representatives is
$(\chi_1 x,\chi_1(s+yw(g)),\chi_1\mu y)$, giving precisely the
restored triangular action on the quotient. Admission follows from
T80B8. Conversely, for an admitted homomorphism $f$ subtract
$\nu(\cdot)f(F)$. The difference $h$ is admitted and $h(F)=0$.
If it were nonzero, let $Y$ be the span of its values, pull back
admission along $Y\hookrightarrow X$, and take a nonzero linear
quotient $Y\to k$. Its composition with $h$ is nonzero since the
values span. It and $\nu\bmod\lambda$ would be independent:
evaluation on $F$ eliminates the coefficient of the degree, and
some value of the other homomorphism eliminates the other coefficient.
This contradicts T80B6. Thus $A_{O(\chi_1)}(X)$ is precisely the
unramified subspace.

Combining the three comparisons gives the native
\texttt{tableLineDegreeTwo.TopCore} for each eligible unramified $K$,
splitting coefficient field and uniformizer, under the residual realization
and the rational period and cohomology inputs used above. Indeed, take
the native $t,h,w,\eta,\mathrm{typ},U,\rho,T_0,\chi$ satisfying
\texttt{TopHyp}. Its three conclusions apply to every finite
$O$-module $X$. For the original lattice, admission is equivalent to
writing the cocycle as a finite sum of integral Bloch--Kato cocycles
tensored with elements of $X$, plus a single coboundary. This is exactly
the tensor image of T80B8, including any supplied lattice whose underlying
submodule is $T_0$ and whose action is $\chi_1^{-1}\rho$. For every
lattice in the line $\chi_c$ with $2\le c\le t$ it is equivalent
to being a coboundary on $G_K$. For every lattice in the first line it
is equivalent to being a coboundary on inertia. A rank-one lattice has an
$O$-basis, and changing that basis transports the entire extension and
all three classes. In each test the jump multiset is
$U_i+[w_i,w_i+t-1]$ and the reference polynomial at $g$ is the
characteristic polynomial of $\rho(g)\oplus\bigoplus_{c=1}^t\chi_c(g)$.
Its role in admission is congruence on the characteristic polynomials
of full lattices, with the multiplicity specified by the native witness.
% END T80B9

% BEGIN T80B10
\paragraph{T80B10. From the native position and polynomial tests to the three extension tests.}
Start with the literal \texttt{TopHyp} used in T80B9. In the unramified
case $K_0=K$ the level-one frame enumerates the $n=[K:\mathbf Q_\ell]$
embeddings. The coefficient field contains their images. Crystalline
comparison identifies each filtered component as a vector space of rank
$r=\dim\rho$, so the specified multiset has $|U_i|=r$ at every letter.
For a true letter filter $U_i$ by $u\le\eta_i$, for a false letter
by $h_i-\eta_i\le u$. Take the complementary filter as cloud. The
ambient condition $1\le u\le h_i-1$ identifies these filters with the
native windows, including their endpoints. Their sum as multisets is
$U_i$. Keep $h,w,\eta,t$ unchanged and use the indicated letter types.
The bounds on the two endpoints put the entire added block in the
short window and give $h_i\ge2$, $\eta_i\ge1$ because $t\ge1$.
Its multiplicity at $a$ is the indicator of $w_i\le a\le w_i+t-1$.
The nonnegative real defining $s_i$ has the same floor whether taken in
$\mathbf N$ or $\mathbf Z$. Hence the three inequalities in
\texttt{PositionHyps} give exactly the size, distance and residue
hypotheses of T80B1. Give the lattice record its original action and
$\chi_1$, and take the integral cyclotomic character $\epsilon$.
The identity $\chi_c=\chi_1\epsilon^{-(c-1)}$ for $1\le c\le t$
follows over $O$ from its stated rational identity by injectivity of
$O\hookrightarrow\operatorname{Frac}(O)$. This constructs the same
numerical and Galois records used in T80B2.

For any continuous full rational representation stable lattices exist:
the compact image takes the elements of a fixed basis into a bounded
set. The $O$-span of their orbits lies between two full lattices,
hence is finitely generated and stable. In any stable lattice basis
the characteristic polynomial is integral and its coefficientwise
reduction is the characteristic polynomial on the reduction. Thus
the polynomial congruence for the full representation says that its
reduction and the reduction of $T_0\oplus\bigoplus_{c=1}^tO(\chi_c)$
have identical characteristic polynomials on each $g$.

Here is the Brauer--Nesbitt step over the actual residue field $k$.
For two modules with matching characteristic polynomials, pass first to
their semisimplifications: a composition series gives block upper
triangular matrices and preserves the determinants. Equality on group
elements gives equality on the group algebra. Explicitly, for a finite
list of matrices $A_j$ attached to group elements and formal variables
$x_j$, the formal determinant
$\det(1-z\sum_j x_j A_j)$ is the product, over primitive cyclic words
$w$ in the indices, of $\det(1-z^{|w|}x_w A_w)$. Here $A_w$ and
$x_w$ are the ordered product and scalar monomial. At any fixed degree
the product is finite. To verify the identity, use universal matrices
in characteristic zero and expand the logarithm as
$-\sum_{m\ge1}z^m\operatorname{tr}((\sum x_jA_j)^m)/m$.
The repetitions of each primitive cyclic word of length $b$ have
$b$ cyclic placements, giving exactly the logarithm of its displayed
factor. The original determinant identity has integer coefficients and
therefore specializes in every characteristic. Each $A_w$ represents
a group element; all factors agree for the two modules. Evaluate the
$x_j$ at coefficients of any group-algebra element.

The finite-dimensional image algebra on the direct sum of the two
semisimple modules is semisimple, since its radical kills all its simple
constituents and the action is faithful. Its primitive central
idempotents project onto the separate isotypic components, one simple
module type for each simple Artinian factor. The characteristic
polynomial of each such idempotent determines the integer dimension of
its image, from the multiplicity of its eigenvalue $1$. Dividing by
the dimension of that simple module over $k$ gives its multiplicity.
Every idempotent lifts to the group algebra, so the two lists of
multiplicities agree over $k$, including when $k$ is not a splitting
field. The reverse implication follows from block triangularity.
Because $O\to k$ is onto, submodules, composition series and their
isomorphisms for these reductions are the same over $O$ and $k$.
The argument applies to every stable lattice: all their characteristic
polynomials lift the same rational polynomial. This is the residual
\texttt{SameSS} condition with its native quantifier on lattices.

Finally, a native full witness is a finite list of full rational
representations and two nested stable lattices $L_2\subset L_1$ in
their product. Giving an equivariant isomorphism to $L_1/L_2$ is
equivalent to a surjective equivariant map from $L_1$ with kernel $L_2$.
A supplied isomorphism from an ambient rational representation to the
product transports both lattices simultaneously. Their quotient is
finite because some power of $\lambda$ times $L_1$ lies in $L_2$.
This recovers the internal admission witness in both directions, with a
separate copy of the prescribed jumps and congruence for every product
factor. For the literal upper triangular action, the map on
representatives is $(x,y)\mapsto(x,y\bmod\lambda^a)$. Its kernel
is $0\times\lambda^aO$ and its equivariance is the identity
\[
 (x,y)\longmapsto
  \big(\rho_W(g)x+y\chi_1(g)z(g),\,\chi_1(g)y\big)
                       \pmod{0\times\lambda^aO} .
\]
It is well-defined when $\lambda^aX=0$. A continuous cocycle with
these finite discrete coefficients is locally constant, and
$\chi_1^{-1}\rho_W$ gives exactly the twisted action in its cocycle
identity. T80B8 therefore applies to every large killing exponent in
the original quantifiers. In particular a coboundary changes the
section of the admitted split sum, so the zero-class converse holds
also for a nonzero representative. T80B9 supplies the corresponding
converse on inertia for the first line.
% END T80B10
